\documentclass[11pt]{article}
\newcommand{\DocShort}{Presentaci\'on semanal}
\usepackage{fontspec}
\usepackage[spanish,es-noshorthands]{babel}
\decimalpoint
\usepackage{geometry}
\geometry{letterpaper, top=2.5cm, bottom=2.5cm, left=2.8cm, right=2.8cm}
\usepackage[expansion=false]{microtype}
\usepackage{parskip}
\setlength{\parskip}{6pt}
\usepackage{amsmath,amssymb}
\usepackage{booktabs}
\usepackage{array}
\usepackage{longtable}
\usepackage{caption}
\usepackage[dvipsnames,table]{xcolor}
\usepackage{enumitem}
\setlist[itemize]{topsep=2pt, itemsep=2pt, parsep=0pt}
\setlist[enumerate]{topsep=2pt, itemsep=2pt, parsep=0pt}
\usepackage{fancyhdr}
\usepackage[hidelinks,pdfencoding=auto]{hyperref}
\usepackage{titlesec}
\usepackage{tcolorbox}
\tcbuselibrary{breakable}
\usepackage{pdfpages}

\definecolor{darkblue}{HTML}{1B3A5C}
\definecolor{medblue}{HTML}{2E6DA4}
\definecolor{lightblue}{HTML}{D6E8F7}
\definecolor{teal}{HTML}{1A7A6E}
\definecolor{lightteal}{HTML}{D0EDEA}
\definecolor{lightgrey}{HTML}{F5F5F5}
\definecolor{midgrey}{HTML}{6F6F6F}
\definecolor{warnyellow}{HTML}{FFF3CD}

\titleformat{\section}
{\large\bfseries\color{darkblue}}
{\thesection}{0.75em}{}[{\color{medblue}\titlerule[0.8pt]}]
\titleformat{\subsection}
{\normalsize\bfseries\color{darkblue}}
{\thesubsection}{0.75em}{}
\titleformat{\subsubsection}
{\normalsize\bfseries\color{medblue}}
{\thesubsubsection}{0.75em}{}

\pagestyle{fancy}
\fancyhf{}
\fancyhead[L]{\small\color{midgrey}ICV 442 Acueductos y Alcantarillados}
\fancyhead[R]{\small\color{midgrey}\DocShort}
\fancyfoot[C]{\small\color{midgrey}\thepage}
\renewcommand{\headrulewidth}{0.4pt}
\renewcommand{\footrulewidth}{0pt}

\rowcolors{2}{lightgrey}{white}
\renewcommand{\arraystretch}{1.18}

\newtcolorbox{infobox}{
  breakable,
  colback=lightblue,
  colframe=medblue,
  boxrule=0.8pt,
  arc=3pt,
  left=8pt,
  right=8pt,
  top=7pt,
  bottom=7pt
}

\newtcolorbox{warnbox}{
  breakable,
  colback=warnyellow,
  colframe=orange!70!black,
  boxrule=0.8pt,
  arc=3pt,
  left=8pt,
  right=8pt,
  top=7pt,
  bottom=7pt
}

\newcommand{\doccover}[3]{
  \begin{center}
  {\small\color{midgrey}Pontificia Universidad Cat\'olica Madre y Maestra\\Escuela de Ingenier\'ia Civil y Ambiental}

  \vspace{1.0cm}
  {\color{medblue}\rule{\linewidth}{2pt}}
  \vspace{0.3cm}

  {\Large\bfseries\color{darkblue}ICV-442-T Acueductos y Alcantarillado}\\[0.35em]
  {\LARGE\bfseries\color{darkblue}#1}\\[0.25em]
  {\large\itshape #2}

  \vspace{0.3cm}
  {\color{medblue}\rule{\linewidth}{2pt}}
  \vspace{0.9cm}
  \end{center}

  \begin{center}
  \begin{tabular}{ll}
  \toprule
  \textbf{Asignatura} & ICV-442-T Acueductos y Alcantarillado \\
  \textbf{Profesor} & Ricardo Rafael Hern\'andez Moreira, PhD \\
  \textbf{Periodo acad\'emico} & Mayo--agosto de 2026 \\
  \textbf{Proyecto} & Dise\~no de sistemas de saneamiento para Las Charcas, Azua \\
  \textbf{Versi\'on} & #3 \\
  \bottomrule
  \end{tabular}
  \end{center}

  \vspace{0.4cm}
}


\begin{document}
\doccover{Formato para presentaci\'on semanal de avance}{Gu\'ia operativa para exposiciones de 5 minutos}{24 de mayo de 2026}

\begin{infobox}
Nombre sugerido del archivo:

\url{ICV442_T[Equipo]_PresentacionSemanal_SemXX_[Apellidos]_v1.0.pdf}
\end{infobox}

\section{Prop\'osito}
\begin{itemize}
\item Este formato organiza las presentaciones semanales de avance de los jueves.
\item La presentaci\'on debe informar progreso real, dificultades t\'ecnicas y pr\'oximos pasos inmediatos.
\item El objetivo no es impresionar con dise\~no visual, sino comunicar con claridad y dentro del tiempo asignado.
\end{itemize}

\section{Regla principal: control del tiempo}
\begin{warnbox}
Cada equipo dispone de \textbf{5 minutos exactos}. Si la presentaci\'on excede ese tiempo, el equipo pierde capacidad de cierre, preguntas y retroalimentaci\'on. La recomendaci\'on operativa es ensayar antes y llegar al aula con una versi\'on ya cronometrada entre \textbf{4:30 y 4:45}.
\end{warnbox}

\begin{itemize}
\item No usar m\'as de 4 diapositivas principales.
\item No dedicar tiempo a contexto general que ya toda la clase conoce.
\item Mostrar solo resultados, decisiones y problemas de la semana.
\item Cada integrante debe saber qu\'e decir y en qu\'e momento intervenir.
\end{itemize}

\section{Estructura sugerida de los 5 minutos}
\begin{center}
\begin{tabular}{p{0.18\linewidth} p{0.36\linewidth} p{0.32\linewidth}}
\toprule
\textbf{Tiempo} & \textbf{Contenido} & \textbf{Apoyo visual sugerido} \\
\midrule
0:00--0:30 & Estado general del equipo y porcentaje de avance frente al cronograma interno. & Diapositiva de resumen con sem\'aforo o porcentaje. \\
0:30--2:00 & Avances t\'ecnicos concretos de la semana. & Plano, captura de EPANET/SWMM, tabla de par\'ametros o esquema. \\
2:00--3:30 & Dificultades encontradas y acciones correctivas. & Lista corta de brechas y respuesta prevista. \\
3:30--5:00 & Pr\'oximos pasos inmediatos y solicitudes de apoyo al profesor. & Cronograma actualizado o lista breve de decisiones pendientes. \\
\bottomrule
\end{tabular}
\end{center}

\section{Plantilla m\'inima por diapositiva}
\subsection*{Diapositiva 1: Estado general}
\begin{itemize}
\item Equipo:
\item Semana:
\item Hito del viernes:
\item Avance acumulado estimado:
\item Estado general: \'optimo / aceptable / con dificultades
\end{itemize}

\subsection*{Diapositiva 2: Avances t\'ecnicos}
\begin{itemize}
\item Decisi\'on o resultado principal 1
\item Decisi\'on o resultado principal 2
\item Evidencia visual o tabla breve
\end{itemize}

\subsection*{Diapositiva 3: Dificultades}
\begin{itemize}
\item Problema detectado
\item Impacto en el proyecto
\item Acci\'on correctiva propuesta
\end{itemize}

\subsection*{Diapositiva 4: Pr\'oximos pasos}
\begin{itemize}
\item Entrega del viernes
\item Tareas internas antes de la pr\'oxima semana
\item Pregunta concreta para el profesor
\end{itemize}

\section{Lista r\'apida de control antes de exponer}
\begin{itemize}
\item La versi\'on presentada coincide con el estado real del trabajo.
\item El equipo puede explicar cualquier valor, criterio o figura mostrada.
\item Las diapositivas son legibles a distancia.
\item El archivo abre sin depender de internet.
\item El ensayo previo dur\'o menos de 5 minutos.
\item Hay una persona controlando el tiempo durante la exposici\'on.
\end{itemize}

\section{Errores que deben evitarse}
\begin{itemize}
\item Leer texto completo de las diapositivas.
\item Mostrar tablas extensas imposibles de interpretar en clase.
\item Gastar 2 o 3 minutos explicando contexto ya conocido.
\item Improvisar el cierre y quedarse sin tiempo para las dificultades o preguntas.
\item Llevar muchas diapositivas y pasar r\'apido por todas sin explicar nada.
\end{itemize}

\section{Nota de uso}
\begin{infobox}
La presentaci\'on semanal no sustituye la bit\'acora ni la entrega del viernes. Debe ser consistente con ambas y servir como insumo para la retroalimentaci\'on inmediata del profesor.
\end{infobox}

\end{document}
